\chapter{Subtipagem e Featherweight Java}

\section{Introdução}

Sistemas de tipos baseados em igualdade exata de tipos são corretos,
mas às vezes rejeitam programas perfeitamente seguros. Se uma função
espera um valor do tipo \(\mathit{Number}\) e fornecemos um
\(\mathit{Int}\), por que o compilador deveria reclamar? Afinal, todo
inteiro é um número.

A \textbf{subtipagem} é o mecanismo que torna esse raciocínio
preciso. Uma relação de subtipagem \(S <: T\) afirma que \emph{todo
  valor do tipo \(S\) pode ser usado com segurança onde um valor do
tipo \(T\) é esperado}: o \textbf{princípio de substituição de
Liskov}. Essa ideia é central em linguagens orientadas a objetos e
permite que subclasses herdem e estendam comportamentos de
superclasses sem quebrar contratos de tipo.

Neste capítulo, introduzimos a subtipagem em duas etapas. Primeiro,
definimos um \textbf{$\lambda$-cálculo com \(\mathsf{Int}\),
\(\mathsf{Bool}\) e funções}, onde \(\mathsf{Bool} <: \mathsf{Int}\),
e estudamos as propriedades de progresso, preservação e
decidibilidade. Em seguida, apresentamos o \textbf{Featherweight Java
(FJ)}, um núcleo formal mínimo de Java que captura a essência da
subtipagem por herança de classes. Finalizamos com a implementação
Haskell completa do verificador de tipos e do interpretador de FJ.

\section{$\lambda$-cálculo com Subtipagem}

\subsection{Sintaxe}

O $\lambda$-cálculo com \(\mathsf{Int}\), \(\mathsf{Bool}\) e funções
tem a seguinte gramática:

\[
  \begin{array}{rcll}
    t & ::= & x & \text{variável} \\
      & |   & n & \text{literal inteiro}\ (n \in \mathbb{Z}) \\
      & |   & \mathsf{true} \mid \mathsf{false} & \text{literais booleanos} \\
      & |   & t + t & \text{adição} \\
      & |   & \mathsf{if}\ t\ \mathsf{then}\ t\ \mathsf{else}\ t & \text{condicional} \\
      & |   & \lambda x {:} T.\, t & \text{abstração tipada} \\
      & |   & t\, t & \text{aplicação}
  \end{array}
\]

Os tipos são:

\[
  \begin{array}{rcll}
    T & ::= & \mathsf{Int}  & \text{inteiros} \\
      & |   & \mathsf{Bool} & \text{booleanos} \\
      & |   & T \to T       & \text{funções}
  \end{array}
\]

Os valores (formas canônicas) são os termos totalmente avaliados:

\[
  \begin{array}{rcll}
    v & ::= & n & \text{inteiro} \\
      & |   & \mathsf{true} \mid \mathsf{false} & \text{booleano} \\
      & |   & \lambda x {:} T.\, t & \text{abstração}
  \end{array}
\]

\subsection{A Relação de Subtipagem}

A subtipagem é definida pela menor relação reflexiva e transitiva que
satisfaz as seguintes regras:

\[ \infer[S-Refl]{T <: T}{} \qquad \infer[S-Trans]{S <: T}{S <: U \quad U <: T}
\]

A regra central do nosso exemplo é que \(\mathsf{Bool}\) é um subtipo
de \(\mathsf{Int}\). Interpretamos \(\mathsf{true}\) como \(1\) e
\(\mathsf{false}\) como \(0\), de modo que todo booleano pode ser
promovido a inteiro sem perda de informação:

\[ \infer[S-Bool]{\mathsf{Bool} <: \mathsf{Int}}{} \]

Para tipos função, a subtipagem é \textbf{contravariante} no domínio
e \textbf{covariante} no contradomínio:

\[ \infer[S-Arrow]{S_1 \to S_2 <: T_1 \to T_2}{T_1 <: S_1 \quad S_2 <: T_2} \]

A intuição para tipos funcionais é a seguinte: se \(f : S_1 \to S_2\)
e queremos usar \(f\) onde se espera \(T_1 \to T_2\), então:

\begin{itemize}
  \item
    O chamador fornecerá argumentos do tipo \(T_1\). Como \(T_1 <:
    S_1\), esses argumentos são aceitos por \(f\).
  \item
    \(f\) retorna \(S_2\). Como \(S_2 <: T_2\), o resultado pode ser
    usado onde se espera \(T_2\).
\end{itemize}

\subsection{Sistema de Tipos}

O contexto de tipagem \(\Gamma\) mapeia variáveis a tipos:
\(\Gamma = x_1 {:} T_1, \ldots, x_n {:} T_n\).

A regra de \textbf{subsunção} é o elo entre a relação de subtipagem e
o sistema de tipos:

\[ \infer[T-Sub]{\Gamma \vdash t : T}{\Gamma \vdash t : S \quad S <: T} \]

Esta regra diz: se \(t\) tem tipo \(S\) e \(S\) é um subtipo de
\(T\), então \(t\) também tem tipo \(T\). As demais regras são:

\[ \infer[T-Var]{\Gamma \vdash x : T}{x {:} T \in \Gamma} \qquad
\infer[T-Int]{\Gamma \vdash n : \mathsf{Int}}{} \]

\[ \infer[T-True]{\Gamma \vdash \mathsf{true} : \mathsf{Bool}}{} \qquad
\infer[T-False]{\Gamma \vdash \mathsf{false} : \mathsf{Bool}}{} \]

\[ \infer[T-Add]{\Gamma \vdash t_1 + t_2 : \mathsf{Int}}{\Gamma \vdash t_1 :
\mathsf{Int} \quad \Gamma \vdash t_2 : \mathsf{Int}} \]

\[ \infer[T-If]{\Gamma \vdash \mathsf{if}\ t_1\ \mathsf{then}\ t_2\
  \mathsf{else}\ t_3 : T}{\Gamma \vdash t_1 : \mathsf{Bool} \quad \Gamma \vdash
t_2 : T \quad \Gamma \vdash t_3 : T} \]

\[ \infer[T-Abs]{\Gamma \vdash \lambda x {:} T_1.\, t : T_1 \to T_2}{\Gamma,\,
  x {:} T_1 \vdash t : T_2} \qquad \infer[T-App]{\Gamma \vdash t_1\, t_2 :
T_2}{\Gamma \vdash t_1 : T_1 \to T_2 \quad \Gamma \vdash t_2 : T_1} \]

Observe que a regra para if exige que os dois ramos tenham o
\textbf{mesmo} tipo \(T\). Graças à subsunção, se os ramos tiverem
tipos distintos mas relacionados por subtipagem, podemos elevá-los ao
supertipo comum antes de aplicar a regra.

\subsection{Semântica Operacional Small-Step}

A semântica procede por redução call-by-value. As regras de
congruência propagam a avaliação para os subtermos:

\[ \infer[E-AddL]{t_1 + t_2 \to t_1' + t_2}{t_1 \to t_1'} \qquad
  \infer[E-AddR]{n + t \to n + t'}{t \to t'} \qquad \infer[E-Add]{n_1 + n_2 \to
n_1 \oplus n_2}{} \]

\[ \infer[E-IfCond]{\mathsf{if}\: t_1\:\mathsf{then}\: t_2\:
    \mathsf{else}\: t_3 \to \mathsf{if}\: t_1'\: \mathsf{then}\: t_2\:
\mathsf{else}\: t_3}{t_1 \to t_1'} \]

\[ \infer[E-IfTrue]{\mathsf{if}\: \mathsf{true}\: \mathsf{then}\: t_2\:
  \mathsf{else}\: t_3 \to t_2}{} \qquad \infer[E-IfFalse]{\mathsf{if}\:
\mathsf{false}\: \mathsf{then}\: t_2\: \mathsf{else}\: t_3 \to t_3}{} \]

\[ \infer[E-AppFun]{t_1\, t_2 \to t_1'\, t_2}{t_1 \to t_1'} \qquad
  \infer[E-AppArg]{v\, t \to v\, t'}{t \to t'} \qquad \infer[E-Beta]{(\lambda x
{:} T.\, t)\, v \to t[x \mapsto v]}{} \]

A subtipagem \(\mathsf{Bool} <: \mathsf{Int}\) implica que booleanos
podem aparecer como operandos de \(+\). Neste caso, o interpretador
os converte numericamente (\(\mathsf{true} \mapsto 1\),
\(\mathsf{false} \mapsto 0\)) antes de efetuar a adição.

\subsection{Propriedades do Sistema de Tipos}

\subsubsection{Lema das Formas Canônicas}

\textbf{Lema (Formas Canônicas).} Seja \(v\) um valor bem tipado em
contexto vazio.

\begin{enumerate}
  \item
    Se \(\vdash v : \mathsf{Int}\), então \(v\) é um literal inteiro
    \(n\) \emph{ou} \(v \in \{\mathsf{true}, \mathsf{false}\}\).
  \item
    Se \(\vdash v : \mathsf{Bool}\) \emph{e} \(v\) não pode ter um
    tipo estritamente menor que \(\mathsf{Bool}\), então
    \(v \in \{\mathsf{true}, \mathsf{false}\}\).
  \item
    Se \(\vdash v : T_1 \to T_2\), então \(v = \lambda x {:} T.\, t\)
    com \(T <: T_1\).
\end{enumerate}

\subsubsection{Progresso}

\textbf{Teorema (Progresso).} Se \(\vdash t : T\), então ou \(t\) é
um valor, ou existe \(t'\) tal que \(t \to t'\).

\emph{Demonstração (esboço).} Por indução na derivação de \(\vdash t : T\).

\begin{itemize}
  \item
    \textbf{T-Var}: impossível em contexto vazio.
  \item
    \textbf{T-Int, T-True, T-False}: \(t\) já é um valor.
  \item
    \textbf{T-Add}: pela hipótese de indução, \(t_1\) e \(t_2\) são
    valores ou reduzem. Se ambos são valores, o Lema das Formas
    Canônicas garante que são representáveis como
    inteiros, então \textbf{E-Add} se aplica.
  \item
    \textbf{T-If}: pelo Lema das Formas Canônicas, se \(t_1\) é um
    valor do tipo \(\mathsf{Bool}\), é \(\mathsf{true}\) ou
    \(\mathsf{false}\), e \textbf{E-IfTrue} ou \textbf{E-IfFalse} se aplicam.
  \item
    \textbf{T-Abs}: \(t\) já é um valor.
  \item
    \textbf{T-App}: se \(t_1\) é um valor do tipo \(T_1 \to T_2\),
    pelo Lema das Formas Canônicas é uma abstração e \textbf{E-Beta} se aplica.
  \item
    \textbf{T-Sub}: segue das hipóteses de indução sobre o subtermo. \(\square\)
\end{itemize}

\subsubsection{Preservação}

\textbf{Lema (Substituição).} Se \(\Gamma, x {:} S \vdash t : T\) e
\(\Gamma \vdash s : S\), então \(\Gamma
\vdash t[x \mapsto s] : T\).

\textbf{Teorema (Preservação).} Se \(\vdash t : T\) e \(t \to t'\),
então \(\vdash t' : T\).

\emph{Demonstração (esboço).} Por indução simultânea na derivação de
tipagem e na redução.

\begin{itemize}
  \item
    \textbf{E-Beta}: a redução \((\lambda x {:} T_1.\, t)\, v \to t[x
    \mapsto v]\) requer \(\vdash v : T_1\). Pelo Lema da
    Substituição, \(\vdash t[x \mapsto v] : T_2\), que é o tipo retornado por \textbf{T-App}.
  \item
    \textbf{E-Add}: \(n_1 + n_2\) reduz ao inteiro resultante, que
    tem tipo \(\mathsf{Int}\) por \textbf{T-Int}.
  \item
    \textbf{E-IfTrue/E-IfFalse}: o branch selecionado tem o tipo
    exigido pela premissa de \textbf{T-If}.
  \item
    Regras de congruência: seguem diretamente das hipóteses de indução.
  \item
    \textbf{T-Sub}: como \(t : S\) e \(S <: T\), após a redução temos
    \(t' : S\) pela hipótese de indução, e \(t' : T\) por
    \textbf{T-Sub}. \(\square\)
\end{itemize}

\subsubsection{Decidibilidade da Subtipagem}

\textbf{Teorema (Decidibilidade).} A relação \(S <: T\) é decidível
para os tipos do nosso cálculo.

\emph{Demonstração.} O conjunto de tipos é indutivamente definido por
três construtores (\(\mathsf{Int}\), \(\mathsf{Bool}\), \(\to\)).
Podemos verificar \(S <: T\) por indução estrutural simultânea em \(S\) e \(T\):

\begin{enumerate}
  \item
    Se \(S = T\), retorne \(\mathsf{true}\) por \textbf{S-Refl}.
  \item
    Se \(S = \mathsf{Bool}\) e \(T = \mathsf{Int}\), retorne
    \(\mathsf{true}\) por \textbf{S-Bool}.
  \item
    Se \(S = S_1 \to S_2\) e \(T = T_1 \to T_2\), retorne
    \(T_1 <: S_1 \land S_2 <: T_2\) por \textbf{S-Arrow} (recursão estrutural,
    logo termina).
  \item
    Caso contrário, retorne \(\mathsf{false}\).
\end{enumerate}

A transitividade não precisa ser verificada explicitamente: a única
cadeia possível de comprimento \(> 1\) é
\(\mathsf{Bool} <: \mathsf{Int} <: \mathsf{Int}\), que
já é coberta pelos casos acima. Logo o algoritmo é correto e termina.
\(\square\)

\section{Featherweight Java}

O \textbf{Featherweight Java (FJ)} foi proposto por Igarashi, Pierce
e Wadler (2001) como um núcleo formal mínimo de Java. Ele captura os
mecanismos essenciais da orientação a objetos (classes, herança,
campos, métodos, criação de objetos e \emph{casts}) numa linguagem
pequena o suficiente para análise formal rigorosa.

FJ omite deliberadamente aspectos que aumentariam a complexidade sem
iluminar os princípios centrais: não há atribuição (todas as
variáveis são imutáveis), não há controle de fluxo, não há
interfaces, não há sobrecarga de operadores e não há tipos primitivos
(tudo é objeto). O que permanece é suficiente para demonstrar
progresso, preservação e a correção dos \emph{casts}.

\subsection{Sintaxe}

\subsubsection{Tipos e Nomes}

Em FJ, \emph{todo tipo é um nome de classe}. O universo de tipos é o
conjunto dos nomes das classes declaradas, mais a classe raiz
\(\mathsf{Object}\) que está sempre presente implicitamente.

\[ T, C, D, E ::= \text{nome de classe} \qquad \mathsf{Object}\ \text{é a raiz
implícita} \]

\subsubsection{Declarações de Classe}

Uma declaração de classe tem a forma:

\[ \mathtt{class}\ C\ \mathtt{extends}\ D\ \{\ \overline{T\ f};\ K\
\overline{M}\ \} \]

onde \(\overline{T\ f}\) é uma lista de campos declarados (com
tipos), \(K\) é o construtor e \(\overline{M}\) é uma lista de
métodos. Usamos a notação com barra para listas: \(\bar{f}\) abrevia
\(f_1, \ldots, f_n\).

O \textbf{construtor} deve ter a forma canônica:

\[ C(\overline{D\ g},\, \overline{T\ f})\ \{\ \mathtt{super}(\bar{g});\
\overline{\mathtt{this}.f = f};\ \} \]

onde \(\bar{g}\) são os campos herdados (passados ao
\(\mathtt{super}\)) e \(\bar{f}\) são os campos próprios de \(C\).

Uma \textbf{declaração de método} tem a forma:

\[ T\ m(\overline{T\ x})\ \{\ \mathtt{return}\ e;\ \} \]

A gramática completa das expressões é:

\[
  \begin{array}{rcll}
    e & ::= & x              & \text{variável} \\
      & |   & e.f            & \text{acesso a campo} \\
      & |   & e.m(\bar{e})   & \text{invocação de método} \\
      & |   & \mathtt{new}\ C(\bar{e}) & \text{criação de objeto} \\
      & |   & (C)\ e         & \text{cast}
  \end{array}
\]

A variável especial \(\mathtt{this}\) é tratada como variável comum
no corpo dos métodos.

\subsubsection{Valores}

O único tipo de valor em FJ é um objeto totalmente construído:

\[ v ::= \mathtt{new}\ C(\bar{v}) \]

Um valor é uma instância de classe cujos campos são, recursivamente,
todos valores.

\subsection{Tabela de Classes}

A \textbf{tabela de classes} \(\mathit{CT}\) é um mapeamento finito
de nomes de classe para suas declarações, construída a partir do
programa. Ela está implícita em todas as regras a seguir.

Definimos quatro funções de consulta:

\textbf{\(\mathit{fields}(C)\)}: lista de \emph{todos} os campos
de \(C\), incluindo os herdados, na ordem de declaração da raiz para a folha:

\[ \mathit{fields}(\mathsf{Object}) = \bullet \qquad \infer{\mathit{fields}(C)
  = \bar{D}\ \bar{g},\, \overline{C'\ f}}{\mathit{CT}(C) = \mathtt{class}\ C\
\mathtt{extends}\ D\ \{\ldots\} \quad \mathit{fields}(D) = \bar{D}\ \bar{g}} \]

\textbf{\(\mathit{mtype}(m, C)\)}: assinatura do método \(m\) em
\(C\) (busca na hierarquia):

\[ \infer{\mathit{mtype}(m, C) = \bar{T} \to T}{\mathit{CT}(C) = \ldots\ T\
  m(\overline{T\ x})\ \{\ldots\}\ \ldots} \qquad \infer{\mathit{mtype}(m, C) =
U}{m \notin \mathit{CT}(C) \quad \mathit{mtype}(m, D) = U} \]

\textbf{\(\mathit{mbody}(m, C)\)}: parâmetros e corpo do método
\(m\) em \(C\):

\[ \infer{\mathit{mbody}(m, C) = (\bar{x}, e)}{\mathit{CT}(C) = \ldots\ T\
m(\overline{T\ x})\ \{\ \mathtt{return}\ e;\ \}\ \ldots} \]

\textbf{Subtipagem \(C <: D\)}: fecho reflexivo-transitivo da
relação \(\mathtt{extends}\):

\[ \infer[S-Refl]{C <: C}{} \qquad \infer[S-Trans]{C <: E}{C <: D \quad D <: E}
  \qquad \infer[S-Extends]{C <: D}{\mathit{CT}(C) = \mathtt{class}\ C\
\mathtt{extends}\ D\ \{\ldots\}} \]

Todo tipo é subtipo de \(\mathsf{Object}\) (por S-Trans e S-Extends
aplicados iterativamente até a raiz).

\subsection{Semântica Operacional}

A semântica é small-step call-by-value. As regras de redução são:

\textbf{Acesso a campo (B-Field):}

\[ \infer[B-Field]{(\mathtt{new}\ C(\bar{v})).f_i \to v_i}{\mathit{fields}(C) =
\overline{C\ f}} \]

O valor \(v_i\) é o \(i\)-ésimo argumento do construtor,
correspondendo ao campo \(f_i\) na lista plana de campos.

\textbf{Invocação de método (B-Invk):}

\[ \infer[B-Invk]{(\mathtt{new}\ C(\bar{v})).m(\bar{u}) \to e[\bar{x} \mapsto
      \bar{u},\, \mathtt{this} \mapsto
  \mathtt{new}\ C(\bar{v})]}{\mathit{mbody}(m, C)
= (\bar{x}, e)} \]

\textbf{Cast (B-Cast):}

\[ \infer[B-Cast]{(C)\ (\mathtt{new}\ D(\bar{v})) \to \mathtt{new}\
D(\bar{v})}{D <: C} \]

O \emph{cast} é verificado em tempo de execução: se \(D\) não é
subtipo de \(C\), uma \(\mathsf{ClassCastException}\) é lançada.

\textbf{Regras de congruência:}

\[ \infer[E-Field]{e_0.f \to e_0'.f}{e_0 \to e_0'} \qquad
\infer[E-InvkRecv]{e_0.m(\bar{e}) \to e_0'.m(\bar{e})}{e_0 \to e_0'} \]

\[ \infer[E-InvkArg]{v_0.m(\ldots, e_i, \ldots) \to v_0.m(\ldots, e_i',
  \ldots)}{e_i \to e_i'} \qquad \infer[E-New]{\mathtt{new}\ C(\ldots,
    e_i, \ldots)
\to \mathtt{new}\ C(\ldots, e_i', \ldots)}{e_i \to e_i'} \]

\[ \infer[E-Cast]{(C)\ e \to (C)\ e'}{e \to e'} \]

\subsection{Sistema de Tipos}

\subsubsection{Tipagem de Expressões}

As regras de tipagem usam um contexto \(\Gamma\) que mapeia variáveis
a tipos. A regra de subsunção, análoga à do $\lambda$-cálculo com
subtipagem, é implícita: em qualquer posição onde se espera um tipo
\(T\), aceita-se um subtipo \(S <: T\).

\[ \infer[T-Var]{\Gamma \vdash x : T}{x {:} T \in \Gamma} \]

\[ \infer[T-Field]{\Gamma \vdash e_0.f_i : T_i}{\Gamma \vdash e_0 : C_0 \quad
\mathit{fields}(C_0) = \overline{T\ f}} \]

\[ \infer[T-Invk]{\Gamma \vdash e_0.m(\bar{e}) : T}{\Gamma \vdash e_0 : C_0
    \quad \mathit{mtype}(m, C_0) = \overline{T} \to T \quad \Gamma
    \vdash \bar{e} :
\bar{C} \quad \bar{C} <: \bar{T}} \]

\[ \infer[T-New]{\Gamma \vdash \mathtt{new}\ C(\bar{e}) : C}{\mathit{fields}(C)
    = \overline{T\ f} \quad \Gamma \vdash \bar{e} : \bar{C} \quad \bar{C} <:
\bar{T}} \]

Para os \emph{casts}, FJ distingue três casos que merecem atenção:

\[ \infer[T-UCast]{\Gamma \vdash (C)\ e_0 : C}{\Gamma \vdash e_0 : D \quad D <:
  C} \qquad \infer[T-DCast]{\Gamma \vdash (C)\ e_0 : C}{\Gamma \vdash e_0 : D
\quad C <: D \quad C \neq D} \]

\[ \infer[T-SCast]{\Gamma \vdash (C)\ e_0 : C}{\Gamma \vdash e_0 : D \quad C
\not<: D \quad D \not<: C \quad \text{(aviso estúpido)}} \]

\begin{itemize}
  \item
    \textbf{T-UCast} (\emph{upcast}): \(D <: C\). O cast é
    estritamente seguro em tempo de compilação: um subtipo sempre
    pode ser tratado como supertipo.
  \item
    \textbf{T-DCast} (\emph{downcast}): \(C <: D\) com \(C \neq D\).
    O cast pode falhar em tempo de execução se o objeto real não for
    instância de \(C\). O verificador de tipos aceita, mas o
    interpretador verifica dinamicamente.
  \item
    \textbf{T-SCast} (\emph{stupid cast}): nenhuma relação de
    subtipagem existe entre \(C\) e \(D\). FJ ainda aceita o programa
    (emitindo um aviso), para não rejeitar código que pode ser
    correto em combinação com tipos parametrizados.
\end{itemize}

\subsubsection{Boa Formação de Métodos}

Um método \(T\ m(\overline{T\ x})\{\ \mathtt{return}\ e;\ \}\) na
classe \(C\) é bem formado se:

\[ \infer[T-Method]{T\ m(\overline{T\ x})\{\ \mathtt{return}\ e;\ \}\
  \mathsf{OK}\ \mathsf{em}\ C}{\overline{x {:} T},\, \mathtt{this}
    {:} C \vdash e
: S \quad S <: T \quad \text{override}(m, D, \overline{T} \to T)} \]

A condição \(\text{override}(m, D, \overline{T} \to T)\) exige que,
se \(m\) já é declarado em uma superclasse \(D\) de \(C\), a
assinatura seja idêntica (FJ não suporta contravariância nos
overrides para simplificar a teoria).

\subsubsection{Boa Formação de Classes}

Uma declaração de classe é bem formada se:

\begin{enumerate}
  \item
    A superclasse existe em \(\mathit{CT}\).
  \item
    O construtor aceita exatamente todos os campos (herdados +
    próprios) na ordem correta, chama \(\mathtt{super}\) com os
    campos herdados e inicializa exatamente os campos próprios.
  \item
    Cada método é bem formado.
\end{enumerate}

\subsubsection{Propriedades}

\textbf{Teorema (Progresso para FJ).} Se \(\vdash e : T\) e \(e\) não
é um valor, então existe \(e'\) tal que \(e \to e'\) ou \(e\) lança
uma \(\mathsf{ClassCastException}\).

A ressalva sobre a exceção de cast é essencial: \emph{downcasts}
podem falhar em tempo de execução mesmo que o programa seja bem
tipado. FJ é, portanto, um sistema de tipos \emph{unsound} no sentido
estrito, mas captura com fidelidade o comportamento real de Java.

\textbf{Teorema (Preservação para FJ).} Se \(\vdash e : T\) e \(e \to
e'\), então \(\vdash e' : T'\) com \(T' <: T\).

O tipo do resultado pode ser um \emph{subtipo} do tipo original,
por exemplo, após um downcast bem-sucedido. Esse enfraquecimento é
necessário e esperado.

\section{Implementação}

\subsection{Sintaxe Abstrata}

O módulo \texttt{FJ.Frontend.Syntax.FJSyntax} define
os tipos Haskell correspondentes à gramática de FJ:

\begin{minted}{haskell}
type ClassName  = String
type FieldName  = String
type MethodName = String
type VarName    = String

newtype FJType = FJType { unFJType :: ClassName }
  deriving (Show, Eq, Ord)

data Program = Program
  { progClasses :: [ClassDecl]
  , progMain    :: Expr
  }

data ClassDecl = ClassDecl
  { cdName    :: ClassName
  , cdSuper   :: ClassName
  , cdFields  :: [(FJType, FieldName)]
  , cdCtor    :: Constructor
  , cdMethods :: [MethodDecl]
  }

data Constructor = Constructor
  { ctorClass  :: ClassName
  , ctorParams :: [(FJType, VarName)]
  , ctorSuper  :: [VarName]
  , ctorInits  :: [(FieldName, VarName)]
  }

data MethodDecl = MethodDecl
  { mdRetType :: FJType
  , mdName    :: MethodName
  , mdParams  :: [(FJType, VarName)]
  , mdBody    :: Expr
  }

data Expr
  = EVar    VarName
  | EField  Expr FieldName
  | EInvk   Expr MethodName [Expr]
  | ENew    ClassName [Expr]
  | ECast   ClassName Expr
\end{minted}

\subsection{Tabela de Classes}

O módulo \texttt{FJ.Frontend.ClassTable.ClassTable}
implementa as quatro funções de consulta definidas na semântica formal:

\begin{minted}{haskell}
type ClassTable = Map.Map ClassName ClassDecl

-- Todos os campos de C (herdados + próprios), raiz primeiro.
classFields :: ClassTable -> ClassName -> Maybe [(FJType, FieldName)]
classFields _  "Object" = Just []
classFields ct c = do
    cd <- Map.lookup c ct
    parentFs <- classFields ct (cdSuper cd)
    return (parentFs ++ cdFields cd)

-- Assinatura de m em C, buscando na hierarquia.
mtype :: ClassTable -> MethodName -> ClassName
      -> Maybe ([FJType], FJType)
mtype _  _ "Object" = Nothing
mtype ct m c = do
    cd <- Map.lookup c ct
    case filter (\md -> mdName md == m) (cdMethods cd) of
      (md:_) -> Just (map fst (mdParams md), mdRetType md)
      []     -> mtype ct m (cdSuper cd)

-- Parâmetros e corpo de m em C.
mbody :: ClassTable -> MethodName -> ClassName
      -> Maybe ([VarName], Expr)
mbody _  _ "Object" = Nothing
mbody ct m c = do
    cd <- Map.lookup c ct
    case filter (\md -> mdName md == m) (cdMethods cd) of
      (md:_) -> Just (map snd (mdParams md), mdBody md)
      []     -> mbody ct m (cdSuper cd)

-- Subtipagem: fecho reflexivo-transitivo de extends.
isSubtype :: ClassTable -> ClassName -> ClassName -> Bool
isSubtype _  c d | c == d   = True
isSubtype _  _ "Object"     = True
isSubtype ct c d =
    case Map.lookup c ct of
      Nothing -> False
      Just cd -> isSubtype ct (cdSuper cd) d
\end{minted}

Note que \texttt{isSubtype ct c "Object"} é sempre
\texttt{True} pelo segundo caso, o que reflecte o
fato de que \texttt{Object} é a raiz da hierarquia.

\subsection{Verificador de Tipos}

O verificador segue diretamente as regras formais. Usamos a mônada
\texttt{Except String} para propagar erros de tipo:

\begin{minted}{haskell}
type TcM a = Except String a
type Env   = Map.Map VarName FJType

tcExpr :: ClassTable -> Env -> Expr -> TcM FJType

-- T-Var
tcExpr _ env (EVar x) =
  case Map.lookup x env of
    Just t  -> return t
    Nothing -> throwError $ "Unbound variable: " ++ x

-- T-Field
tcExpr ct env (EField e f) = do
    FJType c <- tcExpr ct env e
    case classFields ct c of
      Nothing -> throwError $ "Unknown class: " ++ c
      Just fs ->
        case lookup f (map (\(t,fn) -> (fn,t)) fs) of
          Just t  -> return t
          Nothing -> throwError $
            "Field '" ++ f ++ "' not found in class " ++ c

-- T-Invk
tcExpr ct env (EInvk e m args) = do
    FJType c <- tcExpr ct env e
    (paramTypes, retType) <-
      maybe (throwError $ "Method '" ++ m ++ "' not found in " ++ c)
            return
            (mtype ct m c)
    argTypes <- mapM (tcExpr ct env) args
    forM_ (zip argTypes paramTypes) $ \(FJType got, FJType expected) ->
      unless (isSubtype ct got expected) $
        throwError $ "Argument type mismatch: '" ++ got ++
                     "' is not a subtype of '" ++ expected ++ "'"
    return retType

-- T-New
tcExpr ct env (ENew c args) = do
    fs <- maybe (throwError $ "Unknown class: " ++ c)
                return
                (classFields ct c)
    let expectedTypes = map fst fs
    argTypes <- mapM (tcExpr ct env) args
    forM_ (zip argTypes expectedTypes) $ \(FJType got, FJType expected) ->
      unless (isSubtype ct got expected) $
        throwError $ "Constructor argument type mismatch"
    return (FJType c)

-- T-UCast / T-DCast / T-SCast
tcExpr ct env (ECast c e) = do
    FJType _ <- tcExpr ct env e
    return (FJType c)
\end{minted}

A regra de cast aceita todos os três casos (upcast, downcast, stupid
cast), delegando a verificação dinâmica ao interpretador, tal como
especificado pelas regras T-UCast, T-DCast e T-SCast.

A verificação de métodos constrói o ambiente com
\texttt{this} e os parâmetros, verifica o corpo e
checa compatibilidade com eventuais overrides:

\begin{minted}{haskell}
tcMethod :: ClassTable -> ClassName -> MethodDecl -> TcM ()
tcMethod ct c md = do
    let env = Map.fromList $
                ("this", FJType c) :
                map (\(t, x) -> (x, t)) (mdParams md)
    bodyType <- tcExpr ct env (mdBody md)
    unless (isSubtype ct (unFJType bodyType) (unFJType (mdRetType md))) $
      throwError $ "Return type mismatch in method '" ++ mdName md ++ "'"
    case Map.lookup c ct >>= \cd -> mtype ct (mdName md) (cdSuper cd) of
      Nothing -> return ()
      Just (superParams, superRet) ->
        unless (map fst (mdParams md) == superParams &&
                mdRetType md == superRet) $
          throwError $ "Override signature mismatch for '" ++ mdName md ++ "'"
\end{minted}

\subsection{Interpretador}

Os únicos valores em FJ são objetos totalmente construídos:

\begin{minted}{haskell}
data FJValue = FJVal ClassName [FJValue]
\end{minted}

A substituição de variáveis percorre recursivamente a expressão:

\begin{minted}{haskell}
subst :: [(VarName, Expr)] -> Expr -> Expr
subst env (EVar x) =
  case lookup x env of
    Just e  -> e
    Nothing -> EVar x
subst env (EField e f)   = EField (subst env e) f
subst env (EInvk e m as) = EInvk (subst env e) m (map (subst env) as)
subst env (ENew c as)    = ENew c (map (subst env) as)
subst env (ECast c e)    = ECast c (subst env e)
\end{minted}

O núcleo do interpretador implementa as regras B-Field, B-Invk e B-Cast:

\begin{minted}{haskell}
eval :: ClassTable -> Expr -> EvalM FJValue

-- B-New
eval ct (ENew c args) = FJVal c <$> mapM (eval ct) args

-- B-Field: localiza o índice do campo na lista plana
eval ct (EField e f) = do
    v@(FJVal c _) <- eval ct e
    fs <- maybe (throwError $ "Unknown class: " ++ c)
                return (classFields ct c)
    case findIndex (\(_, fn) -> fn == f) fs of
      Nothing -> throwError $ "Field '" ++ f ++ "' not found"
      Just i  -> return (nthArg v i)
  where nthArg (FJVal _ vs) i = vs !! i

-- B-Invk: substitui this e parâmetros no corpo do método
eval ct (EInvk e m args) = do
    v@(FJVal c _) <- eval ct e
    argVals <- mapM (eval ct) args
    (params, body) <-
      maybe (throwError $ "Method '" ++ m ++ "' not found")
            return (mbody ct m c)
    let env = ("this", valToExpr v) :
              zip params (map valToExpr argVals)
    eval ct (subst env body)

-- B-Cast: verifica subtipagem em tempo de execução
eval ct (ECast c e) = do
    v@(FJVal d _) <- eval ct e
    unless (isSubtype ct d c) $
      throwError $ "ClassCastException: '" ++ d ++
                   "' is not a subtype of '" ++ c ++ "'"
    return v
\end{minted}

A função \texttt{valToExpr} converte um valor de
volta a uma expressão \texttt{ENew} para que possa
ser substituído em corpos de métodos:

\begin{minted}{haskell}
valToExpr :: FJValue -> Expr
valToExpr (FJVal c args) = ENew c (map valToExpr args)
\end{minted}

\subsection{REPL}

O executável \texttt{fj} oferece uma interface
interativa para carregar classes e avaliar expressões:

\begin{verbatim}
Featherweight Java Interpreter
  :load <file>   load class declarations from a FJ source file
  :classes       list loaded classes
  :debug         toggle debug output
  :q             quit

Type a FJ expression to type-check and evaluate it.
\end{verbatim}

Um exemplo de sessão com o arquivo \texttt{examples/pair.fj}:

\begin{verbatim}
FJ> :load examples/pair.fj
Loaded 3 class(es): A B Pair
FJ> new Pair(new A(), new B()).fst()
type : Object
--> new A()
FJ> (A) new Pair(new A(), new B()).fst()
type : A
--> new A()
FJ> new Pair(new A(), new B()).swap()
type : Pair
--> new Pair(new B(), new A())
\end{verbatim}

O comando \texttt{:debug} ativa a impressão da
expressão após o parsing, o que é útil para inspecionar a estrutura
da árvore sintática antes da avaliação.

\section{Geração de Código Intermediário para FJ}

O módulo \texttt{FJ.Backend.IR.FJCodegen} compila um programa FJ para
a representação intermediária (IR) usada pelos capítulos anteriores.
A função principal é:

\begin{minted}{haskell}
compileFJ :: Program -> Either String IR.Program
\end{minted}

Para cada programa FJ, o compilador gera quatro categorias de funções IR:
\texttt{C\_new} (construtor de cada classe \(C\)),
\texttt{C\_m\_m} (corpo de cada método \(m\) declarado em \(C\)),
\texttt{dispatch\_m} (despacho dinâmico para cada nome de método \(m\)) e
\texttt{main} (expressão principal do programa).

\subsection{Layout de Objetos em Memória}

Todo objeto FJ é representado como um bloco contíguo de palavras no heap:

\[
  \underbrace{\mathtt{class\_tag}}_{\text{offset } 0}
  \mid
  \underbrace{f_0}_{\text{offset } 8}
  \mid
  \underbrace{f_1}_{\text{offset } 16}
  \mid \cdots \mid
  \underbrace{f_n}_{\text{offset } (n{+}1) \cdot 8}
\]

A palavra na posição 0 guarda uma \emph{tag} inteira única atribuída a
cada classe em tempo de compilação (\texttt{Object} recebe a tag 0).
As palavras seguintes guardam os valores dos campos, incluindo os
herdados, na mesma ordem em que \(\mathit{fields}(C)\) os retorna.
Todos os valores são representados como inteiros de 8 bytes na IR
(ponteiros ou escalares).

O ambiente de compilação encapsula as três informações estáticas
necessárias:

\begin{minted}{haskell}
data CgEnv = CgEnv
  { cgClassTable :: ClassTable
  , cgClassIds   :: Map.Map ClassName Int      -- classe → tag inteiro único
  , cgTypeEnv    :: Map.Map VarName ClassName  -- variável local → tipo estático
  }
\end{minted}

A mônada de geração de código acumula funções geradas e um bloco de
instruções em construção:

\begin{minted}{haskell}
data CgState = CgState
  { cgFuncs   :: [IR.FuncDef]
  , cgStmts   :: [IR.Stmt]
  , cgCounter :: Int
  }

type CgM a = StateT CgState (ExceptT String Identity) a
\end{minted}

A operação \texttt{inNewBlock} isola a geração de um bloco de
instruções sem perturbar o bloco externo, o que é essencial para
compilar o corpo de funções aninhadas:

\begin{minted}{haskell}
inNewBlock :: CgM () -> CgM [IR.Stmt]
inNewBlock m = do
  saved <- gets cgStmts
  modify (\st -> st { cgStmts = [] })
  m
  stmts <- gets cgStmts
  modify (\st -> st { cgStmts = saved })
  return stmts
\end{minted}

\subsection{Construtores}

Para cada classe \(C\) com \(n\) campos (incluindo herdados), o
compilador gera a função \texttt{C\_new}:

\begin{minted}{haskell}
compileConstructor :: CgEnv -> ClassDecl -> CgM ()
compileConstructor env cd = do
    let c      = cdName cd
        cid    = fromJust (Map.lookup c (cgClassIds env))
        Just fs = classFields (cgClassTable env) c
        n      = length fs
        params = ["_p" ++ show i | i <- [0 .. n - 1]]
    body <- inNewBlock $ do
        objT <- fresh
        emit $ IR.MOVE (IR.TEMP objT)
                       (IR.CALL (IR.NAME "alloc") [IR.CONST (n + 1)])
        emit $ IR.MOVE (IR.MEM (IR.TEMP objT)) (IR.CONST cid)
        forM_ (zip [0..] params) $ \(i, p) ->
            emit $ IR.MOVE
                     (IR.MEM (IR.BINOP IR.BAdd (IR.TEMP objT)
                                               (IR.CONST ((i + 1) * 8))))
                     (IR.TEMP p)
        emit $ IR.RETURN [IR.TEMP objT]
    emitFunc (IR.FuncDef (c ++ "_new") params (foldStmts body))
\end{minted}

A função aloca \(n+1\) palavras (uma para a tag e \(n\) para os
campos), escreve a tag na palavra 0 e, para cada campo \(i\), escreve
o parâmetro \texttt{\_p}$i$ no offset \((i+1) \cdot 8\).

\subsection{Métodos}

Para um método \(m\) declarado na classe \(C\), o compilador gera a
função \texttt{C\_m\_m}, onde o receptor é passado como primeiro
parâmetro \texttt{\_this}:

\begin{minted}{haskell}
compileMethod :: CgEnv -> ClassName -> MethodDecl -> CgM ()
compileMethod env c md = do
    let m        = mdName md
        pnames   = map snd (mdParams md)
        ptypes   = map (unFJType . fst) (mdParams md)
        typeEnv  = Map.fromList $ ("this", c) : zip pnames ptypes
        env'     = withTypeEnv typeEnv env
        irParams = "_this" : pnames
    body <- inNewBlock $ do
        emit $ IR.MOVE (IR.TEMP "this") (IR.TEMP "_this")
        resT <- compileExpr env' (mdBody md)
        emit $ IR.RETURN [IR.TEMP resT]
    emitFunc (IR.FuncDef (c ++ "_m_" ++ m) irParams (foldStmts body))
\end{minted}

O \texttt{MOVE (TEMP "this") (TEMP "\_this")} expõe o receptor sob o
nome \texttt{this} para que \texttt{EVar "this"} no corpo do método
seja resolvido corretamente.

\subsection{Despacho Dinâmico}

FJ usa despacho dinâmico: o método executado depende do tipo dinâmico
do receptor, não do tipo estático. O compilador gera uma função
\texttt{dispatch\_m} para cada nome de método \(m\), que lê a tag do
receptor e encadeia comparações para rotear a chamada:

\begin{minted}{haskell}
compileDispatch :: CgEnv -> MethodName -> CgM ()
compileDispatch env m = do
    let ct = cgClassTable env
        definers = [md | cd <- Map.elems ct, md <- cdMethods cd, mdName md == m]
    case definers of
        [] -> return ()
        (md : _) -> do
            let arity    = length (mdParams md)
                dpParams = ["_dp" ++ show i | i <- [0 .. arity - 1]]
                irParams = "_recv" : dpParams
                receivers = classesWithMethod ct m
            body <- inNewBlock $ do
                tagT <- fresh
                emit $ IR.MOVE (IR.TEMP tagT) (IR.MEM (IR.TEMP "_recv"))
                buildDispatchChain env m tagT dpParams receivers
            emitFunc (IR.FuncDef ("dispatch_" ++ m) irParams (foldStmts body))
\end{minted}

A cadeia de despacho para cada classe receptora \(C\) compara a tag lida
com a tag de \(C\) e, em caso de igualdade, realiza uma chamada de cauda
à implementação correta do método:

\begin{minted}{haskell}
buildDispatchChain
    :: CgEnv -> MethodName -> IR.Temp -> [IR.Temp] -> [ClassName] -> CgM ()
buildDispatchChain _   _ _    _      []     =
    emit $ IR.RETURN [IR.CONST 0]   -- inalcançável em programas bem tipados
buildDispatchChain env m tagT dpParams (c : cs) = do
    let impl = fromJust (implementorOf ct m c)
        fn   = impl ++ "_m_" ++ m
        callArgs = map IR.TEMP ("_recv" : dpParams)
    okL   <- freshLabel
    nextL <- freshLabel
    emit $ IR.CJUMP (IR.BINOP IR.BEq (IR.TEMP tagT) (IR.CONST cid)) okL nextL
    emit $ IR.LABEL okL
    emit $ IR.RETURN [IR.CALL (IR.NAME fn) callArgs]
    emit $ IR.LABEL nextL
    buildDispatchChain env m tagT dpParams cs
  where
    ct  = cgClassTable env
    cid = fromJust (Map.lookup c (cgClassIds env))
\end{minted}

A função \texttt{implementorOf} percorre a hierarquia de herança para
encontrar a classe que define diretamente o método \(m\) para
objetos de classe \(C\):

\begin{minted}{haskell}
implementorOf :: ClassTable -> MethodName -> ClassName -> Maybe ClassName
implementorOf _  _ "Object" = Nothing
implementorOf ct m c =
    case Map.lookup c ct of
        Nothing -> Nothing
        Just cd ->
            if any (\md -> mdName md == m) (cdMethods cd)
                then Just c
                else implementorOf ct m (cdSuper cd)
\end{minted}

\subsection{Verificação de Cast}

Para a expressão \texttt{(D) e}, o compilador lê a tag do objeto
resultante e verifica se ela pertence ao conjunto de tags de classes
que são subtipos de \(D\):

\begin{minted}{haskell}
compileExpr env (ECast d e) = do
    ptrT <- compileExpr env e
    let ct       = cgClassTable env
        ids      = cgClassIds env
        validIds = [i | (cn, i) <- Map.toList ids, isSubtype ct cn d]
    okL   <- freshLabel
    failL <- freshLabel
    tagT  <- fresh
    emit $ IR.MOVE (IR.TEMP tagT) (IR.MEM (IR.TEMP ptrT))
    emit $ IR.CJUMP (buildOrChain (IR.TEMP tagT) validIds) okL failL
    emit $ IR.LABEL failL
    emit $ IR.EXP (IR.CALL (IR.NAME "print") [IR.CONST (-1)])
    emit $ IR.RETURN [IR.CONST 0]
    emit $ IR.LABEL okL
    return ptrT
\end{minted}

A função auxiliar \texttt{buildOrChain} constrói uma disjunção de
igualdades de tag para a verificação de subtipagem em tempo de execução:

\begin{minted}{haskell}
buildOrChain :: IR.Expr -> [Int] -> IR.Expr
buildOrChain _   []     = IR.CONST 0
buildOrChain tag [i]    = IR.BINOP IR.BEq tag (IR.CONST i)
buildOrChain tag (i:is) =
    IR.BINOP IR.BOr
             (IR.BINOP IR.BEq tag (IR.CONST i))
             (buildOrChain tag is)
\end{minted}

Em caso de falha, o compilador emite \texttt{print(-1)} (sinaliza
\texttt{ClassCastException}) e retorna 0. Em caso de sucesso, retorna
o ponteiro original intacto.

\subsection{Compilação de Expressões}

A função \texttt{compileExpr} compila cada expressão FJ para um
temporário IR que contém o seu valor:

\begin{minted}{haskell}
compileExpr :: CgEnv -> Expr -> CgM IR.Temp

-- EVar: o nome já é um temporário (parâmetros, "this", variáveis locais)
compileExpr _ (EVar x) = return x

-- ENew: chama o construtor
compileExpr env (ENew c args) = do
    argTs <- mapM (compileExpr env) args
    t <- fresh
    emit $ IR.MOVE (IR.TEMP t)
                   (IR.CALL (IR.NAME (c ++ "_new")) (map IR.TEMP argTs))
    return t

-- EField: lê a palavra no offset do campo
compileExpr env (EField e f) = do
    ptrT <- compileExpr env e
    let idx    = fromJust (findIndex (\(_, fn) -> fn == f) fs)
        offset = (idx + 1) * 8
    t <- fresh
    emit $ IR.MOVE (IR.TEMP t)
                   (IR.MEM (IR.BINOP IR.BAdd (IR.TEMP ptrT) (IR.CONST offset)))
    return t

-- EInvk: chama a função de despacho dinâmico
compileExpr env (EInvk e m args) = do
    recvT <- compileExpr env e
    argTs <- mapM (compileExpr env) args
    t <- fresh
    emit $ IR.MOVE (IR.TEMP t)
                   (IR.CALL (IR.NAME ("dispatch_" ++ m))
                             (map IR.TEMP (recvT : argTs)))
    return t
\end{minted}

Para o acesso a campo, o tipo estático do receptor é reconstruído via
\texttt{typeOfExpr}, que percorre o ambiente local e a tabela de
classes para determinar a classe do receptor e, a partir dela, o
índice do campo \(f\) na lista plana \(\mathit{fields}(C)\).

\section{Conclusão}

A subtipagem é o mecanismo que reconcilia rigor formal e flexibilidade
prática em sistemas de tipos orientados a objetos. A relação \(S <: T\),
formalizada como uma pré-ordem compatível com a estrutura dos tipos,
captura o princípio de substituição de Liskov: qualquer contexto que
espera um valor de tipo \(T\) aceita silenciosamente um valor de tipo
\(S\). A regra de subsunção \textbf{T-Sub} é o mecanismo pelo qual essa
compatibilidade é injetada no sistema de tipos, e a contravariância no
domínio de tipos função é a consequência inevitável de se tratar funções
como contratos bidirecionais entre produtor e consumidor.

O Featherweight Java demonstra que esses princípios se aplicam
igualmente à subtipagem por herança de classes. Com apenas cinco formas
de expressão (variável, acesso a campo, invocação de método, criação
de objeto e cast), FJ captura o núcleo formal de Java com precisão
suficiente para demonstrar progresso e preservação. A ressalva do
progresso é reveladora: bem-tipicidade não impede \emph{downcasts}
inválidos de falhar em tempo de execução, o que reflete com fidelidade
o comportamento real de Java. A preservação, por sua vez, enuncia uma
versão enfraquecida da invariante de tipo: o tipo do resultado pode ser
um \emph{subtipo} do tipo original, pois um downcast bem-sucedido
refina a informação disponível sobre o objeto.

A implementação Haskell evidencia a correspondência sistemática entre
especificação formal e código. A tabela de classes é uma função de
consulta que realiza computacionalmente as metafunções
\(\mathit{fields}\), \(\mathit{mtype}\) e \(\mathit{mbody}\) da teoria;
o verificador de tipos é uma transliteração direta das regras de
derivação; e o interpretador implementa as regras de redução B-Field,
B-Invk e B-Cast quase literalmente, com \texttt{subst} realizando a
substituição simultânea de variáveis especificada na regra B-Invk.

O compilador FJ para representação intermediária fecha o ciclo entre
teoria e geração de código. A escolha de representar objetos como blocos
lineares de palavras com uma \emph{tag} de classe na posição zero é
simultaneamente a mais simples e a mais eficiente: acesso a campos é uma
aritmética de ponteiros; despacho dinâmico é uma leitura de tag seguida
de uma cadeia de comparações; e verificação de cast é uma disjunção de
igualdades de tag. Essa representação é, em essência, a mesma usada por
implementações industriais de Java e C++. A função
\texttt{compileFJ} organiza a geração em quatro passagens independentes
(construtores, métodos, funções de despacho e \texttt{main}), o que
reflete a modularidade da própria teoria de FJ e facilita a verificação
de correção de cada fase isoladamente.
